\section{Related Work}
\label{sec:related}
\begin{table}[t]
\centering
\caption{Representative systems expose different organization choices.}
\label{tab:related-work}
\footnotesize
\setlength{\tabcolsep}{2pt}
\renewcommand{\arraystretch}{1.08}
\begin{tabular}{@{}>{\raggedright\arraybackslash}p{.205\columnwidth}>{\raggedright\arraybackslash}p{.14\columnwidth}>{\raggedright\arraybackslash}p{.275\columnwidth}>{\raggedright\arraybackslash}p{.32\columnwidth}@{}}
\toprule
System & Ops./\allowbreak object & Decomposition basis & Organization / resources \\
\midrule
FFTW~\cite{fftw2005} & FFT, DCT/\allowbreak DST & Recursive plans & Codelet composition; empirical planning \\
\addlinespace[2pt]
VkFFT~\cite{vkfft2023,vkfft} & FFT, DCT & Grouped strided transforms & Layout choices \\
\addlinespace[2pt]
GPU-NTT~\cite{gpuntt2023,gpuntt2025} & NTT & Merge / Four-Step & Locality; kernel count; block shape \\
\midrule
TensorIR~\cite{tensorir2023} & Tensor ops. & Block signatures & Outer transformations; tensorization \\
\addlinespace[2pt]
Welder~\cite{welder2023} & DNN graphs & Tile connections; shape propagation & Hierarchical reuse; traffic / footprint \\
\addlinespace[2pt]
Timeloop~\cite{timeloop2019} & DNN layers & Workload / mapping separation & Spatial / temporal tiling; resource costs \\
\midrule
\system{} & FFT, NTT, FWHT & Stage bits; closed regions; intersections & Stage / data reuse; live state, traffic, core / ready work \\
\bottomrule
\end{tabular}
\end{table}

Related systems optimize transform formulas, arithmetic cores, tensor programs,
and hardware mappings.

\noindent\textbf{Transform generation and specialized execution.}
FFTW combines codelets with empirical planning; SPIRAL/FFTX generates
transforms and surrounding computation from algebraic descriptions
\cite{fftw2005,spiral2005,fftx2018}. VkFFT generates GPU implementations with
grouped strided transforms and layout choices. GPU-NTT optimizes locality,
kernel count, block size, and shape
\cite{vkfft2023,vkfft,gpuntt2023,gpuntt2025}.
FlashFFTConv combines tensor-core FFT decomposition, batch/hidden tiling,
and memory-aware fusion~\cite{flashfftconv2024}. HadaCore maps FWHT chunks
to warps and vectors to threadblocks with shared-memory transposes
\cite{hadacore2024}. These systems optimize both arithmetic and organization.

Table~\ref{tab:related-work} separates their representative operators, units,
stage decompositions, independent-work mappings, and boundary mechanisms.
A GPU kernel is a launched program; a tile subgraph groups computation and
values for scheduling and residency. Kernels can execute tile subgraphs, and
larger subgraphs can span kernels. Stage fusion can retain intermediates while
reusing execution resources over time, so it does not alone establish a
spatial pipeline.

\noindent\textbf{Tensor compilation and tile scheduling.}
TVM separates computation from scheduling~\cite{tvm2018}; Ansor searches
hierarchical schedule sketches~\cite{ansor2020}; ROLLER constructs
hardware-aligned tiles~\cite{roller2022}. TensorIR provides block signatures,
outer transformations, and tensorization~\cite{tensorir2023}. Welder connects
producer/consumer tiles and schedules reuse through traffic and footprint
estimates~\cite{welder2023}. These establish tiling and reusable compute
interfaces. Their generality alone does not verify transform-specific support.
For butterflies, the changing pair relation identifies the values needed by
a stage group and the producers that a successor must await, before choosing
a physical tile or arithmetic core. This separates dependency construction
from a particular kernel implementation while preserving the operator's
arithmetic, direction, ordering, and precision contract.

\noindent\textbf{Hierarchical dataflow and positioning.}
Eyeriss analyzes data movement under controlled resources; Timeloop separates
workload, architecture, and hierarchical spatial/temporal mappings
\cite{eyeriss2016,timeloop2019}. Applying these principles to butterflies
requires legal resident tile subgraphs and their remaining communication. GPU
blocks also share resources, precluding an assumption of dedicated pipeline
hardware.

Our distinction concerns the common dependency interface. SPIRAL/FFTX derives
implementations from algebraic transform decompositions, including surrounding
computation~\cite{spiral2005,fftx2018}. Timeloop evaluates hierarchical mappings
from workload and architecture descriptions~\cite{timeloop2019}. These precedents
motivate separating arithmetic from organization; that separation alone is not
our novelty. We derive the butterfly-specific link between dependencies and
resource demand: stage-bit sets construct dependency-closed regions, and their
intersections derive communication and readiness
(Section~\ref{sec:method}). Spatial and temporal organization then determines
the demands evaluated against hardware constraints. The contribution is this butterfly-specific construction, not the general
idea of spatial/temporal dataflow or a claim that prior systems cannot express
equivalent programs. It retains specialized arithmetic across workloads and
hardware constraints. The fixed-core experiments in Section~\ref{sec:core-organization} isolate this organization
effect from arithmetic substitution.
